\scanpage{44}\mspage{43}
This topology is finer than $\tau\deleted{'}$, i.e. if $X$ is closed in $\tau$, it satisfies these conditions : clear. It induces $\tau'_0$, i.e. if $X\subset T$, it satisfies these conditions if and only if it is closed for $\tau'_0$ : clear. Thus $\tau'\in\underline{\mathscr{T}}$. Let us prove that it is a greatest element (coarsest topology). i.e that for \deleted{every} $\tau''\in\underline{\mathscr{T}}$, the closed sets for $\tau'$ are closed sets for $\tau''$. Indeed, we have
\[
X=B\deleted{\cap}\cup\overline{X\capU}^{(\tau)}
\]
$B$ being closed in $T$. (for $\tau'_0$\inserted{\deleted{\unsure{therefore}} for $\tau^{\prime\prime\unsure{,}}\mid\overset{\sim}{T}$}), and $T$ being \uline{closed} in $\deleted{\unsure{T|\tau}}S$ for $\tau$ hence \inserted{a fortiori} for $\tau''$, $B$ is closed in $S$ for $\tau''$.
% slanted
\margin{(This is the first time we use $T$ closed in $S$}
Likewise $\overline{X\cap U}^{(\tau)}$ being closed for $\tau$, is so for $\unsure{:|;}\tau''$, so $X=B\cup\overline{X\cap U}^{(\tau)}$ is so too. This proves a) b) c).

Let us prove d)\inserted{(the beginning already proved)}, recalling that for $x,y\in T$,
\[
\begin{aligned}
x\underset{\tau'}{\leq}y&\Longleftrightarrow\overline{x}^{(\tau')}\subset\overline{y}^{(\tau')}\\
&\Longleftrightarrow x\in\overline{y}^{(\tau')}.
\end{aligned}
\]
Since $\tau'\mid U\ne\tau\mid U$, and $\tau'\mid T=\tau'_0$, the cases $x,y\in T$ and $x,y\in U$ are already clear. There remains the case where either $x\in U,y\in T$, or the reverse. The first case cannot occur, for if $y\in T$, $\overline{y}^{(\tau')}\inT$ ($T$ \uline{closed} in $S$ for $\tau'$) so $x\notin\overline{y}^{(\tau')}$ \deleted{since} if $x\in U$. So there remains the case $[x\in T,]$ $y\in U$, prove $x\underset{\tau'}{\leq}y\Longleftrightarrow x\underset{\tau'}{\leq}y$
\margin{NB this is true without assuming $x\inT$}
